%-----------------------------------------------------------------------------------------------
\makeatletter
\immediate\write18{datelog > \jobname.info} % site script for $(date '+%Y-%m-%d %Hh%Mm%Ss')
\makeatother
%-----------------------------------------------------------------------------------------------
%-----------------------------------------------------------------------------------------------
\usetheme{Copenhagen}
\usepackage{beamercolorthemeCNThermSci2}
\usefonttheme{serif}
%-----------------------------------------------------------------------------------------------

%-----------------------------------------------------------------------------------------------
%-----------------------------------------------------------------------------------------------
\usetheme{Copenhagen}
\usepackage{beamercolorthemeUTF2}
\usefonttheme{serif}
%-----------------------------------------------------------------------------------------------
\usepackage[utf8]{inputenc}
\usepackage[greek,french,english,brazil]{babel} % last becomes the active one
\usepackage{pslatex}
\usepackage{amssymb,amsmath}
\usepackage{soul}
\usepackage[squaren,Gray,cdot]{SIunits}
\usepackage[nice]{nicefrac}
\usepackage{tikz}
\usepackage{amscd}
\usepackage{stmaryrd}
\usepackage{scalerel}
\usepackage{xspace}
%-----------------------------------------------------------------------------------------------


%-----------------------------------------------------------------------------------------------
%-----------------------------------------------------------------------------------------------
% Mathematical
%-----------------------------------------------------------------------------------------------
\newcommand{\vet}[1]{\underline{{#1}}}
\newcommand{\mat}[1]{\underline{\underline{{#1}}}}
\newcommand{\cub}[1]{\underline{\underline{\underline{{#1}}}}}
\newcommand{\eqdef}{{\ensuremath\stackrel{\text{\tiny def}}{=}}}
%-----------------------------------------------------------------------------------------------
% Linguistic
%-----------------------------------------------------------------------------------------------
\newcommand{\GRtxt}[1]{\begin{otherlanguage}{greek}{{#1}}\end{otherlanguage}}
\newcommand{\FRtxt}[1]{\begin{otherlanguage}{french}{{#1}}\end{otherlanguage}}
%-----------------------------------------------------------------------------------------------
% Presentation
%-----------------------------------------------------------------------------------------------
\newcommand{\BkgImgH}[1]{% Places an image centered on the slide background filling the height
    \usebackgroundtemplate{\parbox{\paperwidth}{%
        \vspace*{1sp}\centering\includegraphics[height=\paperheight]{{#1}}
}}}
\newcommand{\BkgImgW}[1]{% Places an image centered on the slide background filling the width
    \usebackgroundtemplate{\parbox{\paperwidth}{%
        \vspace*{1sp}\centering\includegraphics[width=\paperwidth]{{#1}}
}}}
\newcommand{\ArtEndH}[3]{% Transitions to plain image (last) slide: #1:prefix #2,#3:extensions
    \BkgImgH{root/../art/#1.#2}
    \frame<handout:0>[plain]{%
        \transdissolve\vspace*{72mm}\color{white}\scriptsize\bf\input{root/../art/#1.#3}}
    \usebackgroundtemplate{\mbox{~}}
}
\newcommand{\ArtEndW}[3]{% Transitions to plain image (last) slide: #1:prefix #2,#3:extensions
    \BkgImgW{root/../art/#1.#2}
    \frame<handout:0>[plain]{%
        \transdissolve\vspace*{72mm}\color{white}\scriptsize\bf\input{root/../art/#1.#3}}
    \usebackgroundtemplate{\mbox{~}}
}
\newcommand{\ImgColW}[3]{% Inserts a full-width image in a column
    \includegraphics[width=\columnwidth]{root/../art/#1.#2}\\[-0.5\baselineskip]
    \parbox{\columnwidth}{\tiny\hfill\scalebox{0.85}{\input{root/../art/#1.#3}}}
}
\newcommand{\txtpic}[1]{%
    \fcolorbox{lightgray}{white!90!black}{{#1}} 
}
%-----------------------------------------------------------------------------------------------


%-----------------------------------------------------------------------------------------------
\title{D.01.01 -- Fundamentos de Refrigeração}
\subtitle{Refrigeração e Condicionamento de Ar}
\author{Prof.~C.~Naaktgeboren, PhD}
\date{{\scriptsize\tt%
    \includegraphics[height=6.0mm]{cc/by-nc-nd-88x31.pdf}\\[\smallskipamount]
    https://github.com/CNThermSci/ApplThermSci\\
    Compiled on \input{\jobname.info}
}}
%-----------------------------------------------------------------------------------------------
\begin{document}
%-----------------------------------------------------------------------------------------------
\logo{%
    \parbox{158mm}{% There's a 1mm gap on each side of the 160mm x 90mm slide logo line
        \mode<beamer>{
            \includegraphics[height=6.0mm]{root/00-res/UTFPR/UTFPR-logo-D.pdf}\hfill%
            \includegraphics[height=9.0mm]{root/00-res/logo/CNThermSci-logo-A.pdf}%
        }
        \mode<handout>{
            \includegraphics[height=6.0mm]{root/00-res/UTFPR/UTFPR-logo-W.pdf}\hfill%
            \includegraphics[height=9.0mm]{root/00-res/logo/CNThermSci-logo-W.pdf}%
        }
    }
} % The (delineated, alpha), or washed-out logos
%-----------------------------------------------------------------------------------------------
\frame{\titlepage}
%-----------------------------------------------------------------------------------------------

%-----------------------------------------------------------------------------------------------
\frame{\tableofcontents}
%-----------------------------------------------------------------------------------------------

%-----------------------------------------------------------------------------------------------
\frame{%
    \hfill\par
    % !bib LST 2>/dev/null | paste -s -d',' | sed 's|,|, |g' | j -i8 96 
    Esta apresentação baseia-se primordialmente na referência~\cite{2016-FentonDL-ASHRAE},
    \alert{Capítulo 1} (tópico de leitura).
    \hfill
}
%-----------------------------------------------------------------------------------------------

%-----------------------------------------------------------------------------------------------
\section{Sistemas e Processos de Refrigeração}
%-----------------------------------------------------------------------------------------------

%-----------------------------------------------------------------------------------------------
\subsection{Introdução}
%-----------------------------------------------------------------------------------------------

    % !j 96 -i8
    %-------------------------------------------------------------------------------------------
    \begin{frame}{Introdução à Refrigeração}\vspace*{-2em}
        \begin{itemize}
            \item<1-> Definições;
            \item<1-> Item 2;
            \item<1-> Item 3.
        \end{itemize}
    \end{frame}
    %-------------------------------------------------------------------------------------------

%-----------------------------------------------------------------------------------------------
\subsection{Classificação dos Sistemas}
%-----------------------------------------------------------------------------------------------

    % !j 96 -i8
    %-------------------------------------------------------------------------------------------
    \begin{frame}[allowframebreaks]{Sistemas por Compressão de Vapor -- }\vspace*{-2em}
        \begin{itemize}
            \item<1-> Definições;
            \item<1-> Item 2;
            \item<1-> Item 3.
        \end{itemize}
    \end{frame}
    %-------------------------------------------------------------------------------------------

    % !j 96 -i8
    %-------------------------------------------------------------------------------------------
    \begin{frame}[allowframebreaks]{Sistemas a Ar (Gás) -- }\vspace*{-2em}
        \begin{itemize}
            \item<1-> Definições;
            \item<1-> Item 2;
            \item<1-> Item 3.
        \end{itemize}
    \end{frame}
    %-------------------------------------------------------------------------------------------

    % !j 96 -i8
    %-------------------------------------------------------------------------------------------
    \begin{frame}[allowframebreaks]{Sistemas de Absorção -- }\vspace*{-2em}
        \begin{itemize}
            \item<1-> Definições;
            \item<1-> Item 2;
            \item<1-> Item 3.
        \end{itemize}
    \end{frame}
    %-------------------------------------------------------------------------------------------

    % !j 96 -i8
    %-------------------------------------------------------------------------------------------
    \begin{frame}{Sistemas Termo-Elétricos}\vspace*{-2em}
        \begin{itemize}
            \item<1-> Definições;
            \item<1-> Item 2;
            \item<1-> Item 3.
        \end{itemize}
    \end{frame}
    %-------------------------------------------------------------------------------------------

    % !j 96 -i8
    %-------------------------------------------------------------------------------------------
    \begin{frame}{Resfriadores Evaporativos}\vspace*{-2em}
        \begin{columns}
        \column{0.55\textwidth}
            Coluna com 55\% de largura:\\[\smallskipamount]
            \begin{itemize}
                \item<1-> Item 1;
                \item<2-> Item 2;
                \item<3-> Item 3.
            \end{itemize}
        \column{0.45\textwidth}
            Coluna com 45\% de largura.
            \begin{center}
                \begin{figure}
                    \includegraphics[width=0.5\textwidth]{www/moringa-01.jpg}
                \end{figure}
            \end{center}
        \end{columns}
    \end{frame}
    %-------------------------------------------------------------------------------------------

%-----------------------------------------------------------------------------------------------
\subsection{Processos de Refrigeração}
%-----------------------------------------------------------------------------------------------

    % !j 96 -i8
    %-------------------------------------------------------------------------------------------
    \begin{frame}{Ciclo de Refrigeração por Compressão de Vapor}\vspace*{-2em}
        \begin{itemize}
            \item<1-> Definições;
            \item<1-> Item 2;
            \item<1-> Item 3.
        \end{itemize}
    \end{frame}
    %-------------------------------------------------------------------------------------------

    % !j 96 -i8
    %-------------------------------------------------------------------------------------------
    \begin{frame}{Ciclo Bomba de Calor}\vspace*{-2em}
        \begin{itemize}
            \item<1-> Definições;
            \item<1-> Item 2;
            \item<1-> Item 3.
        \end{itemize}
    \end{frame}
    %-------------------------------------------------------------------------------------------

%-----------------------------------------------------------------------------------------------
\section{Aplicações de Refrigeração}
%-----------------------------------------------------------------------------------------------

%-----------------------------------------------------------------------------------------------
\subsection{Classificação por Capacidade}
%-----------------------------------------------------------------------------------------------

    % !j 96 -i8
    %-------------------------------------------------------------------------------------------
    \begin{frame}{O Ramo de Refrigeração}\vspace*{-2em}
        \begin{itemize}
            \item<1-> Doméstico, \alert{menos de $\unit{20}{\kilo\watt}$};
            \item<2-> Comercial, \alert{mais de $\unit{20}{\kilo\watt}$};
            \item<3-> Industrial, \alert{de pequeno a muito grande}.
        \end{itemize}
    \end{frame}
    %-------------------------------------------------------------------------------------------

%-----------------------------------------------------------------------------------------------
\subsection{Classificação por Aplicação}
%-----------------------------------------------------------------------------------------------

    % !j 96 -i8
    %-------------------------------------------------------------------------------------------
    \begin{frame}{Aplicações de Refrigeração}\vspace*{-2em}
        \begin{itemize}
            \item<1-> Condicionamento de ar residencial;
            \item<2-> Condicionamento de ar veicular;
            \item<3-> Condicionamento de ar de médios e grandes edificações;
            \item<4-> Transporte de cargas;
            \item<5-> Refrigeração residencial;
            \item<6-> Refrigeração de máquinas de venda;
            \item<7-> Ar-condicionado industrial;
            \item<8-> Refrigeração industrial;
        \end{itemize}
    \end{frame}
    %-------------------------------------------------------------------------------------------

    % !j 96 -i8
    %-------------------------------------------------------------------------------------------
    \begin{frame}{Ar-condicionado Industrial}\vspace*{-2em}
        \begin{itemize}
            \item<1-> Item 1;
            \item<2-> Item 2;
            \item<3-> Item 3.
        \end{itemize}
    \end{frame}
    %-------------------------------------------------------------------------------------------

    % !j 96 -i8
    %-------------------------------------------------------------------------------------------
    \begin{frame}{Refrigeração Industrial}\vspace*{-2em}
        \begin{itemize}
            \item<1-> Item 1;
            \item<2-> Item 2;
            \item<3-> Item 3.
        \end{itemize}
    \end{frame}
    %-------------------------------------------------------------------------------------------

%-----------------------------------------------------------------------------------------------
\section{Referências e Tópicos de Leitura}
%-----------------------------------------------------------------------------------------------

    %------------------------------------------------------------------------------------------
    \begin{frame}[allowframebreaks]{Referências -- }
        \bibliographystyle{unsrt}
        \setbeamertemplate{bibliography item}{\insertbiblabel}
        \bibliography{bibfile.bib}
    \end{frame}
    %------------------------------------------------------------------------------------------

    % Finishes with stunning image, with credit
    \ArtEndWSE{pexels-emiliano-arano-1295138}{jpg}{txt}

%-----------------------------------------------------------------------------------------------
\end{document}
%-----------------------------------------------------------------------------------------------

